\subsection{The Reduced Norm is a Polynomial Function}

Lemma 3. --- Let L be an extension of K, and let I be a set and \(\mathbf{T} = (\mathrm{T}_i)_{i\in \mathrm{I}}\) a family of variables. We have \(\mathrm{K}(\mathbf{T})\cap \mathrm{L}[\mathbf{T}] = \mathrm{K}[\mathbf{T}]\) .

Let P and Q be elements of K{[}T{]} with \(Q \neq 0\) . The coefficients of the polynomials R in L{[}T{]} such that P = QR are the solutions of a system of linear equations with coefficients in K. Consequently, if there exists a polynomial \(R \in L[T]\) such that P = QR, then there also exists one in K{[}T{]} (II, §8, No.~5, p.~321, Proposition 6). This proves the inclusion \(\mathrm{K}(\mathbf{T}) \cap \mathrm{L}[\mathbf{T}] \subset \mathrm{K}[\mathbf{T}]\) ; the reverse inclusion is obvious.

Recall that the reduced characteristic polynomial of an element a of A can be written as

\[
\mathrm{Pcrd} _ {\mathrm{A/K}} (a; \mathrm{X}) = \sum_ {r = 0} ^ {n} (- 1) ^ {r} b _ {r} (a) \mathrm{X} ^ {n - r}\tag{46}
\]

and that we have

\[
b _ {0} (a) = 1, \quad b _ {1} (a) = \mathrm{Trd} _ {\mathrm{A/K}} (a), \quad b _ {n} (a) = \mathrm{Nrd} _ {\mathrm{A/K}} (a).
\]

PROPOSITION 6. --- For every integer r such that \(0 \leqslant r \leqslant n\) , the mapping \(b_{r}\) from A to K is a homogeneous polynomial mapping of degree r. In particular, the reduced norm is a homogeneous polynomial mapping of degree n from A to K.

Let \((e_{i})_{i\in\mathrm{I}}\) be a basis of A over K and \(\mathbf{T}=(\mathrm{T}_{i})_{i\in\mathrm{I}}\) a family of variables.

Lemma 4. --- Let u be the element \(\sum_{i\in I}T_{i}\otimes e_{i}\) of the central simple \(\mathrm{K}(\mathbf{T})\) -algebra \(\mathrm{A}_{(\mathrm{K}(\mathbf{T}))}\) . Let P be the reduced characteristic polynomial of u. Then P belongs to the ring \(\mathrm{K}[\mathbf{T}][\mathrm{X}]\) ; viewed as an element of the ring \(\mathrm{K}[\mathbf{T},\mathrm{X}]\) , it is homogeneous of degree n.

We choose an extension L of K and an L-algebra isomorphism \(\theta\) from \(\mathrm{A}_{(\mathrm{L})}\) to \(\mathbf{M}_{n}(\mathrm{L})\) . We denote by \(\overline{\theta}: \mathrm{A}_{(\mathrm{L}(\mathbf{T}))} \to \mathbf{M}_{n}(\mathrm{L}(\mathbf{T}))\) the isomorphism of \(\mathrm{L}(\mathbf{T})\) -algebras deduced from \(\theta\) by extension of scalars. By Corollary 1 of VIII, p.~342, we have

\[
\mathrm{P} (\mathrm{X}) = \chi_ {\overline {{\theta}} (u)} (\mathrm{X}) = \det (\mathrm{X} I _ {n} - \overline {{\theta}} (u)) = \det \left(\mathrm{X} I _ {n} - \sum_ {i \in \mathrm{I}} \mathrm{T} _ {i} \theta (1 \otimes e _ {i})\right).\tag{47}
\]

Since the matrices \(\theta(1\otimes e_{i})\) belong to \(\mathbf{M}_{n}(\mathrm{L})\) , this formula shows that P is a homogeneous polynomial of degree n in L{[}T,X{]}. It also belongs to \(\mathrm{K}(\mathbf{T})[\mathrm{X}]\) and can be written as \(\mathrm{P}(\mathrm{X})=\sum_{j\geqslant0}c_{j}\mathrm{X}^{j}\) , where each \(c_{j}\) belongs to the intersection \(\mathrm{K}(\mathbf{T})\cap\mathrm{L}[\mathbf{T}]\) . By Lemma 3, each of the elements \(c_{j}\) belongs to K{[}T{]}; Lemma 4 follows.

Lemma 5. --- For every extension \(K'\) of K and every element \((t_{i})_{i\in I}\) of \(K^{I}\) , we have

\[
\mathrm{Pcrd} _ {\mathrm{A} _ {(\mathrm{K} ^ {\prime})} / \mathrm{K} ^ {\prime}} \Bigl (\sum_ {i \in \mathrm{I}} t _ {i} \otimes e _ {i} \Bigr) = \mathrm{P} ((t _ {i}) _ {i \in \mathrm{I}}, \mathrm{X}).\tag{48}
\]

Let \(\varphi : \mathrm{K}[\mathbf{T}] \to \mathrm{K}'\) be the unique K-algebra homomorphism that sends \(\mathrm{T}_i\) to \(t_i\) for every \(i \in \mathrm{I}\) ; it defines on \(\mathrm{K}'\) the structure of a \(\mathrm{K}[\mathbf{T}]\) -algebra. The \(\mathrm{K}'\) -algebra \(\mathrm{A}_{(\mathrm{K}[\mathbf{T}])(\mathrm{K}')}\) can be identified with \(\mathrm{A}_{(\mathrm{K}')}\) (transitivity of extension of scalars), where the element \(1 \otimes (\sum \mathrm{T}_i \otimes e_i)\) is identified with the element \(\sum t_i \otimes e_i\) of \(\mathrm{A}_{(\mathrm{K}')}\) . We denote by \(\overline{\varphi}: \mathrm{K}[\mathbf{T}][\mathrm{X}] \to \mathrm{K}'[\mathrm{X}]\) the K-algebra homomorphism deduced from \(\varphi\) . By formula (21) of III, §9, No.~3, p.~544, the characteristic polynomial of \(\sum t_i \otimes e_i\) with respect to the \(\mathrm{K}'\) -algebra \(\mathrm{A}_{(\mathrm{K}')}\) is the image by \(\overline{\varphi}\) of the characteristic polynomial of \(\sum \mathrm{T}_i \otimes e_i\) with respect to the \(\mathrm{K}[\mathbf{T}]\) -algebra \(\mathrm{A}_{(\mathrm{K}[\mathbf{T}])}\) , that is, of \(\mathrm{P}^n\) . In other words, we have

\[
\mathrm{Pc} _ {\mathrm{A} _ {(\mathrm{K} ^ {\prime})} / \mathrm{K} ^ {\prime}} \left(\sum_ {i \in \mathrm{I}} t _ {i} \otimes e _ {i}; \mathrm{X}\right) = \mathrm{P} ((t _ {i}) _ {i \in \mathrm{I}}, \mathrm{X}) ^ {n};\tag{49}
\]

Lemma 5 then follows from Lemma 2 of VIII, p.~339.

Consider the specific case \(K' = K\) of Lemma 5. We have

\[
\mathrm{Pcrd} _ {\mathrm{A/K}} \Bigl (\sum_ {i \in \mathrm{I}} t _ {i} e _ {i}; \mathrm{X} \Bigr) = \mathrm{P} ((t _ {i}) _ {i \in \mathrm{I}}, \mathrm{X})\tag{50}
\]

for every element \((t_{i})_{i\in\mathrm{I}}\) of \(K^{I}\) . Since the polynomial P in \(K[T,X]\) is homogeneous of degree n, it can be expanded uniquely as

\[
\mathrm{P} (\mathbf {T}, \mathrm{X}) = \sum_ {r = 0} ^ {n} (- 1) ^ {r} \mathrm{B} _ {r} (\mathbf {T}) \mathrm{X} ^ {n - r},\tag{51}
\]

where \(B_{r}\) is a homogeneous polynomial of degree r in K{[}T{]}. By formulas (46), (50), and (51), we have

\[
b _ {r} \Bigl (\sum_ {i \in \mathrm{I}} t _ {i} e _ {i} \Bigr) = \mathrm{B} _ {r} ((t _ {i}) _ {i \in \mathrm{I}})
\]

for every element \((t_i)_{i\in \mathrm{I}}\) of \(\mathbf{K}^{\mathrm{I}}\) . Proposition 6 follows.

Remark. --- Let \(K'\) be a commutative K-algebra. Every element t of \(\mathrm{A}_{(\mathrm{K}^{\prime})}\) can be written as \(\sum_{i\in I}t_{i}\otimes e_{i}\) , where \((t_{i})\in\mathrm{K}^{\prime\mathrm{I}}\) . It follows from the proof of Lemma 5 that the characteristic polynomial \(\mathrm{Pc}_{\mathrm{A}_{(\mathrm{K}^{\prime})}/\mathrm{K}^{\prime}}(t;\mathrm{X})\) is equal to \(\mathrm{P}((t_{i}),\mathrm{X})^{n}\) .
% semantic-fragments: legacy-input-seam
\relax{}
